\documentclass[11pt]{article}
\usepackage[margin=1in]{geometry}
\usepackage{amsmath,amssymb,amsthm}
\usepackage[T1]{fontenc}
\usepackage{lmodern}
\usepackage{microtype}
\usepackage{booktabs}
\usepackage[colorlinks=true,linkcolor=blue,citecolor=blue,urlcolor=blue]{hyperref}
\hypersetup{pdftitle={Graceful 4-coloring of planar graphs is NP-complete},pdfauthor={P. Graham}}
\newcommand{\GFC}{\textup{\textsc{Graceful 4-Colorability}}}
\newtheorem{theorem}{Theorem}
\newtheorem{proposition}[theorem]{Proposition}
\newtheorem{lemma}[theorem]{Lemma}
\theoremstyle{remark}
\newtheorem{remark}[theorem]{Remark}
\title{Graceful 4-coloring of planar graphs is NP-complete}
\author{P. Graham\\ \small\url{https://orcid.org/0009-0007-8493-1505}}
\date{November 3, 2026}
\begin{document}
\maketitle
\begin{abstract}
A graceful $k$-coloring of a graph is a proper vertex coloring with colors
$1,\dots,k$ such that, at every vertex, the incident edges receive distinct
differences of their endpoint colors. Cyriac Antony, Laavanya D. and Devi
Yamini S. proved that deciding graceful $k$-colorability is NP-complete for
planar graphs when $k\ge 5$, and for $k=4$ on graphs of maximum degree three,
and they left open the complexity of graceful 4-colorability of planar graphs.
We prove that this problem is NP-complete, even for planar graphs of maximum
degree three (inputs, as in their formulation, may be disconnected). The
hardness comes from the reduction of \'Edouard Bonnet, Dibyayan Chakraborty
and Julien Duron, who proved that deciding whether a cubic planar graph has a
perfect matching cut (a perfect matching that is exactly the set of edges
leaving some vertex set) is NP-hard; we apply their construction to every
instance. What we add is an exact correspondence: for every finite cubic graph
$H$, the graph obtained by subdividing each edge of $H$ twice has a graceful
4-coloring if and only if $H$ has a perfect matching cut. This equivalence is
proved in Lean from the definitions.

This work has not yet been reviewed by experts.
The proofs, the Lean formalization and the text were produced by AI models (from OpenAI and Anthropic), working under the author's direction.
\end{abstract}
\medskip
\noindent\textit{2020 Mathematics Subject Classification.} 05C15, 05C78, 68Q17.
\textit{Keywords.} Graceful coloring, perfect matching cut, planar graphs, NP-completeness.

\section{Introduction}

All graphs are finite and simple. For a graph $G$ and a map
$f\colon V(G)\to\{1,2,3,4\}$, the \emph{difference} of an edge $uv$ is
$|f(u)-f(v)|$. The map $f$ is a \emph{graceful $4$-coloring} if adjacent
vertices receive distinct colors and any two edges sharing an endpoint have
distinct differences. Thus a vertex of color $1$ with neighbors colored
$2,3,4$ is fine, while a vertex of color $2$ cannot have three neighbors,
because only the differences $1$ and $2$ are available to it. Colors need
not be distinct on all of $V(G)$, and not all four need be used. The
problem \GFC{} asks whether a given graph has a
graceful $4$-coloring. Graceful colorings were introduced by Bi, Byers,
English, Laforge and Zhang~\cite{Bi}. The definition above is the informal
definition of Antony, Laavanya and Devi Yamini: a proper coloring whose edge
differences form a proper edge coloring \cite[Section~1]{ALD}. Their formal
definition can be stronger on graphs with triangles, but the two agree on the
triangle-free graphs used here, so Theorem~\ref{thm:main} holds for both
(Remark~\ref{rem:ald}).

A \emph{perfect matching cut} of a graph $H$ is a perfect matching $M$ of $H$
for which there is a set $A\subseteq V(H)$ such that $M$ is exactly the set
of edges with one endpoint in $A$ and one outside $A$. Bonnet, Chakraborty
and Duron call this ``a perfect matching that is also a cutset''
\cite[Section~2]{BCD}. The \emph{double subdivision} $S_2(H)$ replaces each
edge $uv$ of $H$ by a path $u\,x\,y\,v$ through two new vertices.

\subsection{Main results}

\begin{theorem}\label{thm:main}
\GFC{} is NP-complete on planar graphs. It remains
NP-complete on planar graphs of maximum degree at most three.
\end{theorem}

We use the construction of Bonnet, Chakraborty and Duron~\cite{BCD} without
their splitting step, which is a Turing reduction. This keeps our reduction
many-one, but the graphs it produces may be disconnected. This matches the
problem as posed, whose input is an arbitrary graph; the connected case is not
treated. The proof rests on the following exact equivalence, which holds for
all cubic graphs, planar or not.

\begin{theorem}\label{thm:equiv}
Let $H$ be a finite cubic graph. Then $S_2(H)$
has a graceful $4$-coloring if and only if $H$ has a perfect matching cut.
\end{theorem}

For example, $K_4$ has no perfect matching cut (every edge cut of $K_4$ has
zero, three or four edges, while a perfect matching has two), so $S_2(K_4)$
has no graceful $4$-coloring. The triangular prism (two triangles joined by
three edges) has one, those three joining edges, so its double subdivision has
a graceful $4$-coloring.

\subsection{Where the question comes from}

Antony, Laavanya and Devi Yamini proved that ``Graceful 5-colorablity [sic] is
NP-complete for 3-regular 3-connected planar bipartite graphs'' and that
``For $k\ge 6$, Graceful $k$-colorablity is NP-complete for planar bipartite
graphs of maximum degree $k-2$'' \cite[Theorems~4 and~5]{ALD}. For four
colors they proved that ``Graceful 4-colorablity is NP-complete for the class
of 2-degenerate graphs of maximum degree 3'' \cite[Theorem~6]{ALD}, by a
reduction from Positive Not-All-Equal 3-Sat E4~\cite{DD}; the graphs this
reduction produces need not be planar. Their
abstract ends: ``The complexity of checking whether a planar graph is graceful
4-colorable remains open.'' The same sentence closes the abstract of the
published conference version by Cyriac Antony, Jacob Antony, Laavanya D.\ and
Devi Yamini S.~\cite{AALD}. Theorem~\ref{thm:main} answers this question.

Pantoja, Dantas and Luiz also proved hardness results for graceful colorings
of subcubic graphs; neither concerns planar graphs with four colors. Pantoja's
thesis, which presents their three papers \cite{PDL371,PDL374,PDL381} among
its main contributions \cite{Pantoja}, attributes to \cite{PDL374} the
statement ``The G5CP is NP-complete even when restricted to 3-connected
bipartite cubic planar graphs'' \cite[Theorem~10]{Pantoja}, where G5CP is the
graceful 5-coloring problem, and to \cite{PDL381} the statement ``The
\textup{\textsc{Graceful 4-Coloring}} problem is NP-complete, even for graphs
with maximum degree 3'' \cite[Theorem~34]{Pantoja}, together with a
polynomial-time algorithm for trees of maximum degree three
\cite[Theorem~33]{Pantoja}. Other recent work on graceful colorings concerns
trees \cite{BNS,Popa}, corona graphs \cite{LDcorona} and a class of bipartite
graphs built by an operation called unfolding \cite{LDN}.

\subsection{What is new}

The input we build on is the Bonnet--Chakraborty--Duron theorem
\cite[Theorem~1]{BCD}: ``Perfect Matching Cut is NP-hard in 3-connected cubic
bipartite planar graphs.'' Their reduction supplies all of the hardness. Our contribution
is Theorem~\ref{thm:equiv}, the exact correspondence between graceful
$4$-colorings of $S_2(H)$ and perfect matching cuts of $H$. Double subdivision
keeps planarity and keeps the maximum degree at three, so the hardness transfers. The step that needs care is the forward direction: a graceful
$4$-coloring of $S_2(H)$ must be shown to produce a perfect matching that is
an entire cut and not merely contained in one. This is a counting argument on
the differences at the original vertices (Section~\ref{sec:equiv}).

We use their construction without the extra hypothesis they make for
$3$-connectivity, which we do not need (Section~\ref{sec:hard}).

\section{Proof of Theorem~\ref{thm:equiv}}\label{sec:equiv}

Let $H$ be cubic and let $f$ be a graceful $4$-coloring of $S_2(H)$. Call the
vertices of $H$ \emph{original}. For an original vertex $u$ and an edge $uv$
of $H$, the \emph{port} of $u$ toward $v$ is the edge from $u$ to the first new
vertex on the path replacing $uv$, and its difference is the port difference.

\begin{lemma}\label{lem:star}
In a graceful $4$-coloring, a vertex with three neighbors has color $1$ or
$4$, its three incident edges have differences $1,2,3$, and its neighbors
receive the three other colors.
\end{lemma}

\begin{proof}
A vertex of color $c$ has at most $|\{|c-d|: d\neq c,\ 1\le d\le 4\}|$ distinct
positive differences available. This is $3$ for $c\in\{1,4\}$ and $2$ for
$c\in\{2,3\}$. Three neighbors need three distinct differences, so
$c\in\{1,4\}$, and then the differences $1,2,3$ are realized by the three
other colors.
\end{proof}

The first claim of Lemma~\ref{lem:star} is the case $\Delta=3$ of
\cite[Lemma~6]{Pantoja}, attributed there to \cite{PDL381}; the same counting
appears in \cite[proof of Construction~1]{ALD}.

\begin{lemma}\label{lem:path}
Let $u\,x\,y\,v$ be a path whose interior vertices have degree two, colored
by a graceful $4$-coloring with $f(u),f(v)\in\{1,4\}$. Up to the symmetry
$c\mapsto 5-c$ and reversing the path, the possibilities are exactly those in
Table~\ref{tab:path}.
\end{lemma}

\begin{table}[h]
\centering
\begin{tabular}{ccccc}
\toprule
type & $f(u),f(v)$ & $f(x),f(y)$ & differences along the path & port differences \\
\midrule
I & $1,1$ & $2,4$ & $1,2,3$ & $1$ and $3$ \\
II & $1,1$ & $3,4$ & $2,1,3$ & $2$ and $3$ \\
III & $1,4$ & $3,2$ & $2,1,2$ & $2$ and $2$ \\
\bottomrule
\end{tabular}
\caption{Graceful colorings of a twice subdivided edge with extreme endpoint colors.}
\label{tab:path}
\end{table}

\begin{proof}
At $x$ the two differences $|f(u)-f(x)|$ and $|f(x)-f(y)|$ must differ, and
similarly at $y$, and consecutive colors must differ. If $f(u)=f(v)=1$, then
$f(x),f(y)\in\{2,3,4\}$ are distinct; $(2,3)$ fails at $x$ (differences
$1,1$) and $(3,2)$ fails at $y$, leaving $(2,4)$, $(3,4)$ and their reversals.
If $f(u)=1$ and $f(v)=4$, then $f(x)=4$ would give $y$ two neighbors of color
$4$, and $f(y)=1$ is excluded in the same way, so $f(x),f(y)\in\{2,3\}$ are
distinct; $(2,3)$ gives differences $1,1,1$, and $(3,2)$ remains.
\end{proof}

\begin{proof}[Proof of Theorem~\ref{thm:equiv}, forward direction]
By Lemma~\ref{lem:star} every original vertex has color $1$ or $4$ and exactly
one port of each difference $1,2,3$. By Lemma~\ref{lem:path}, a port of
difference $1$ lies on a path of type I in Table~\ref{tab:path}, whose
other port has difference $3$. This pairs each of the $|V(H)|$ ports of
difference $1$ with a distinct port of difference $3$. There are also exactly
$|V(H)|$ ports of difference $3$, so all of them are used by paths of type I,
and no path of type II occurs. Hence every port of difference $2$ lies on a
path of type III. Let $M$ be the set of edges of $H$ whose paths are of
type III. Each original vertex has exactly one port of
difference $2$, so $M$ is a perfect matching. Let $A$ be the set of original
vertices of color $1$. Paths of types I and II join endpoints of equal color
and paths of type III join endpoints of different colors, so $M$
is exactly the set of edges leaving $A$.
\end{proof}

\begin{proof}[Converse direction]
Let $M$ be a perfect matching of $H$ equal to the set of edges leaving $A$.
Removing $M$ leaves a graph in which every vertex has degree two, a disjoint
union of cycles, and each cycle lies inside $A$ or inside its complement.
Orient each cycle cyclically. Give the vertices of $A$ color $1$ and the others
color $4$. Color the path replacing a cycle edge $u\to v$ by $1,4,2,1$ inside
$A$ and $4,1,3,4$ outside $A$; color the path replacing a matching edge from
$u\in A$ to $v\notin A$ by $1,3,2,4$. Along these paths the differences are
$3,2,1$ and $2,1,2$, so each new vertex sees two distinct differences, and
consecutive colors differ. At an original vertex the outgoing cycle port has
difference $3$, the incoming cycle port $1$ and the matching port $2$. This is
a graceful $4$-coloring of $S_2(H)$.
\end{proof}

\section{Proof of Theorem~\ref{thm:main}}\label{sec:hard}

\paragraph{Membership in NP.} A certificate is a color in $\{1,2,3,4\}$ for
each vertex. Checking that adjacent vertices have distinct colors and that
edges at each vertex have distinct differences takes time polynomial in the
size of the graph.

\paragraph{Double subdivision.}

\begin{lemma}\label{lem:s2}
Let $H$ be a graph with $n$ vertices and $m$ edges.
\begin{enumerate}
\item[(a)] If $H$ is planar, then $S_2(H)$ is planar.
\item[(b)] $S_2(H)$ has $n+2m$ vertices and $3m$ edges, every new vertex has
degree two, and every original vertex keeps its degree. From adjacency lists,
$S_2(H)$ is computed in linear time.
\end{enumerate}
\end{lemma}

\begin{proof}
For (a), draw $H$ in the plane, choose two interior points on each edge arc
and declare them the new vertices. Part (b) is immediate from the
construction.
\end{proof}

\paragraph{The cited hardness theorem.} We use \cite[Theorem~1]{BCD}
through its proof, which is a reduction from Monotone Not-All-Equal 3SAT-E4.
An instance is a set of clauses, each consisting of three distinct variables,
with each variable in exactly four clauses; it is \emph{positive} if some
truth assignment gives every clause at least one true and at least one false
variable. Bonnet, Chakraborty and Duron cite its NP-completeness from Darmann
and Döcker~\cite{DD} \cite[Section~3.2]{BCD}. For an instance $I$ they
construct ``in polynomial time'' a graph $G(I)$ \cite[Section~3.2]{BCD}. The
proof of \cite[Lemma~13]{BCD} shows that $G(I)$ is cubic, planar and
bipartite, and \cite[Lemmas~26 and~28]{BCD} show that $G(I)$ has a perfect
matching cut if and only if $I$ is positive. These two lemmas rest on the
preceding results of \cite[Section~3]{BCD}; of these, \cite[Lemma~7]{BCD} is
stated for cubic bipartite graphs and \cite[Lemma~3]{BCD} for plane graphs, so
the equivalence uses the planarity and bipartiteness of $G(I)$ from the proof
of \cite[Lemma~13]{BCD}.
Bonnet, Chakraborty and Duron also assume that no variable is a cut vertex of
the variable--clause incidence graph of $I$, justified by a splitting step
that they call ``a polynomial-time Turing reduction'' \cite[Section~3.2]{BCD}.
That hypothesis is used only in the proof of \cite[Lemma~12]{BCD}, that
$G(I)$ is $3$-connected, which the proof of their Lemma~13 quotes for
$3$-connectivity alone. The cubicity, planarity and bipartiteness arguments do
not cite their Lemma~12, and no other lemma of theirs uses the hypothesis. The bipartiteness argument shows that
every bounded face has even length, counting along exposed paths of the
gadgets. The proof of their Lemma~12 records two facts that do not depend on the
hypothesis: every gadget is $2$-connected, and adjacent gadgets are joined by
two disjoint connector edges. Removing one vertex therefore leaves every gadget
connected and keeps a connector edge between any two adjacent gadgets, so each
component of $G(I)$ is $2$-connected, the faces of its drawing are bounded by
cycles, and that count shows that each component, and so $G(I)$, is bipartite.
(Numbers of lemmas in \cite{BCD} refer to arXiv:2302.11667v1.)

\begin{proposition}[from the proof of {\cite[Theorem~1]{BCD}}]\label{prop:bcd}
For every instance $J$ of Monotone Not-All-Equal 3SAT-E4, the graph $G(J)$ is
cubic and planar, it is computed from $J$ in polynomial time, and it has a
perfect matching cut if and only if $J$ is positive.
\end{proposition}

\begin{proof}
The construction and its running time are in \cite[Section~3.2]{BCD},
cubicity and planarity are shown in the proof of \cite[Lemma~13]{BCD}, and
the equivalence is \cite[Lemmas~26 and~28]{BCD}, which use the bipartiteness
shown in the same proof.
\end{proof}

\begin{proof}[Proof of Theorem~\ref{thm:main}]
Given $J$, compute $G(J)$ by Proposition~\ref{prop:bcd} and output
$S_2(G(J))$. Since $G(J)$ is cubic, Proposition~\ref{prop:bcd} and
Theorem~\ref{thm:equiv} give
\begin{align*}
 J\text{ is positive}
 &\iff G(J)\text{ has a perfect matching cut}\\
 &\iff S_2(G(J))\text{ has a graceful 4-coloring}.
\end{align*}
By Lemma~\ref{lem:s2}, $S_2(G(J))$ is planar, has maximum degree three and is
computed in polynomial time. Together with membership in NP, this proves both
statements of Theorem~\ref{thm:main}.
\end{proof}

\begin{remark}\label{rem:ald}
Antony, Laavanya and Devi Yamini also give a formal version of the definition:
the coloring is injective on each closed neighborhood $N[v]$, and the
difference labeling is injective on all edges of $G[N[v]]$
\cite[Section~1]{ALD}. On graphs with triangles it can be stronger than the
definition used here, since it also separates two edges inside $N[v]$ that
share no endpoint. On triangle-free graphs the two agree: the edges of
$G[N[v]]$ are then the edges at $v$, and distinct differences at $v$ force
distinct colors on its neighbors. Every double subdivision is triangle-free,
so the reduction above proves NP-hardness for both versions. Membership in NP
holds for both, because each condition is checked on the closed neighborhoods
in polynomial time. Hence Theorem~\ref{thm:main} holds under either of their
formulations.
\end{remark}

\section{Formal verification}\label{sec:formal}

The Lean project accompanying this paper uses Lean~4 \cite{Lean} and Mathlib
\cite{Mathlib}. Its main theorem is
\begin{verbatim}
theorem Antony24.graceful4_s2_iff_pmc {V : Type*} [Finite V]
    (G : SimpleGraph V) (hG : Cubic G) :
    Graceful4 (S2 G) <-> PMC G
\end{verbatim}
Here \texttt{SimpleGraph V} is Mathlib's simple graph on the vertex type
\texttt{V}; the only assumption on \texttt{V} is finiteness, so empty and
disconnected graphs are included. \texttt{Cubic G} says each vertex has exactly
three distinct neighbors. \texttt{Graceful4 G} says some map to a four-element
palette is proper on edges and gives distinct differences to two distinct
neighbors of each vertex; a proved lemma shows this is the same as labels
$1,\dots,4$. \texttt{S2 G} has vertex type \texttt{V} plus the ordered pairs
$(u,v)$ of adjacent vertices, with the edge $uv$ becoming the path
$u\,(u,v)\,(v,u)\,v$. \texttt{PMC G} says some two-sided split of the vertices
gives each vertex exactly one neighbor on the other side; a proved lemma shows
this is a Mathlib perfect matching equal to the whole cut of a vertex set. So
the theorem is Theorem~\ref{thm:equiv}, and the project also states it with
these unfolded forms on both sides.

The project also proves the counts and degrees of Lemma~\ref{lem:s2}(b) and
that every double subdivision is triangle-free, as used in
Remark~\ref{rem:ald}; the theorem names are listed in the project's README.
The build checks that the main theorems depend only on the axioms
\texttt{propext}, \texttt{Classical.choice} and \texttt{Quot.sound}, and
checks the two examples of Section~1 by kernel evaluation.

Outside Lean are \cite[Theorem~1]{BCD} and its proof, including its use of
\cite{DD}; Proposition~\ref{prop:bcd}; the planarity and running-time parts of
Lemma~\ref{lem:s2}; membership in NP; and the notion of NP-completeness
itself.

To reproduce the check, install Lean through elan and run
\texttt{lake exe cache get} and \texttt{lake build} in the \texttt{lean}
directory. With the downloaded Mathlib cache the build takes under a minute and
about 5\,GB of memory.

\begin{thebibliography}{99}
\bibitem{AALD} Cyriac Antony, Jacob Antony, Laavanya D.\ and Devi Yamini S.,
Graceful coloring is computationally hard, in: Algorithms and Discrete Applied
Mathematics (CALDAM 2025), Lecture Notes in Computer Science 15536, Springer, 2025, 13--24,
\href{https://doi.org/10.1007/978-3-031-83438-7_2}{doi:10.1007/978-3-031-83438-7\_2}.
\bibitem{ALD} Cyriac Antony, Laavanya D.\ and Devi Yamini S., Graceful coloring
is computationally hard, \href{https://arxiv.org/abs/2407.02179v2}{arXiv:2407.02179v2}, 2024.
\bibitem{BNS} Vitaly A. Baransky, Ilia A. Nasyrov and Tatiana A. Senchonok,
Locks and forbidden configurations in 4-graceful trees, Ural Math.\ J.\ 12
(2026), no.~1, 26--44,
\href{https://doi.org/10.15826/umj.2026.1.003}{doi:10.15826/umj.2026.1.003}.
\bibitem{Bi} Z. Bi, A. Byers, S. English, E. Laforge and P. Zhang, Graceful colorings of graphs, J. Combin. Math. Combin. Comput.\
101 (2017), 101--119.
\bibitem{BCD} \'Edouard Bonnet, Dibyayan Chakraborty and Julien Duron, Cutting
Barnette graphs perfectly is hard, Theoret.\ Comput.\ Sci.\ 1010 (2024), 114701,
\href{https://doi.org/10.1016/j.tcs.2024.114701}{doi:10.1016/j.tcs.2024.114701};
\href{https://arxiv.org/abs/2302.11667v1}{arXiv:2302.11667v1}.
\bibitem{DD} Andreas Darmann and Janosch Döcker, On a simple hard variant of
Not-All-Equal 3-Sat, Theoret.\ Comput.\ Sci.\ 815 (2020), 147--152,
\href{https://doi.org/10.1016/j.tcs.2020.02.010}{doi:10.1016/j.tcs.2020.02.010}.
\bibitem{LDcorona} Laavanya D.\ and Devi Yamini S., Graceful coloring of some
corona graphs: an algorithmic approach, Commun.\ Comb.\ Optim.\ 11 (2026),
no.~2, 315--334,
\href{https://doi.org/10.22049/cco.2024.29155.1866}{doi:10.22049/cco.2024.29155.1866}.
\bibitem{LDN} Laavanya D., Devi Yamini S.\ and Muhammad Faisal Nadeem, Graceful
colorings of a subclass of bipartite graphs, in: Nonlinear Complex Systems,
Computational Intelligence and Their Applications, Springer Proceedings in
Mathematics \& Statistics, Springer, 2026, 237--247,
\href{https://doi.org/10.1007/978-981-95-5891-9_18}{doi:10.1007/978-981-95-5891-9\_18}.
\bibitem{Mathlib} The mathlib Community, The Lean mathematical library, in:
Proceedings of the 9th ACM SIGPLAN International Conference on Certified
Programs and Proofs (CPP 2020), 367--381.
\bibitem{Lean} Leonardo de Moura and Sebastian Ullrich, The Lean~4 theorem
prover and programming language, in: Automated Deduction (CADE 28), Lecture
Notes in Computer Science 12699, Springer, 2021, 625--635.
\bibitem{PDL371} Paola T. Pantoja, Simone Dantas and Atílio G. Luiz, Bounds on
graceful $k$-colorings of graphs, Discrete Appl.\ Math.\ 371 (2025), 218--231,
\href{https://doi.org/10.1016/j.dam.2025.04.032}{doi:10.1016/j.dam.2025.04.032}.
\bibitem{PDL374} Paola T. Pantoja, Simone Dantas and Atílio G. Luiz, Graceful
colorings of graphs with maximum degree three, Discrete Appl.\ Math.\ 374
(2025), 75--96,
\href{https://doi.org/10.1016/j.dam.2025.05.013}{doi:10.1016/j.dam.2025.05.013}.
\bibitem{PDL381} Paola T. Pantoja, Simone Dantas and Atílio G. Luiz, On graceful
colorings in subcubic trees, Discrete Appl.\ Math.\ 381 (2026), 98--111,
\href{https://doi.org/10.1016/j.dam.2025.11.027}{doi:10.1016/j.dam.2025.11.027}.
\bibitem{Pantoja} Paola Tatiana Pantoja Huaynoca, Graceful colorings of graphs:
structure, bounds, and complexity, doctoral thesis, Universidade Federal
Fluminense, Niterói, 2026,
\href{https://sipos.app/pgmat/administracao/defesas/documentos/2026.03.23-14.56.04Tese\%20Paola\%20Tatiana\%20Pantoja\%20H.pdf}{open copy at sipos.app}.
\bibitem{Popa} Alexandru Popa and Andrei Popa, A polynomial time algorithm for
computing the graceful coloring number on trees, SSRN preprint 7181532, 2026,
\href{https://doi.org/10.2139/ssrn.7181532}{doi:10.2139/ssrn.7181532}.
\end{thebibliography}
\end{document}
